% Generated by make_book.py from Lecture5/lecture5.tex; edit the lecture, not this file.
\chapter{Differential Equations}\label{L5-chapter}
% The group keeps settings such as \emergencystretch inside the chapter.
\begingroup
\section*{Introduction}
\addcontentsline{toc}{section}{Introduction}

Many quantities in machine learning change continuously in time, or can be understood as if they did: the parameters of a network trained by gradient descent with small steps, the hidden state of a residual network as it passes through many layers, the samples of a generative model as they are transformed from noise into data, and the states of a sampling algorithm as it explores a distribution. Differential equations describe such dynamics. In this lecture:
\begin{itemize}
    \item we solve ordinary differential equations (ODEs) $\mathbf{x}' = F(\mathbf{x})$, in particular linear systems, with the eigenvalues and eigenvectors of Lecture~1;
    \item we see that gradient descent is Euler's method for the gradient flow $\mathbf{x}' = -\nabla f(\mathbf{x})$, which explains its step size limit;
    \item we study how the flow of a vector field changes volumes and probability densities, using the divergence of Lecture~3; this is the basis of continuous normalizing flows;
    \item we add noise: Brownian motion, stochastic differential equations (SDEs) and the Ornstein--Uhlenbeck process, and Langevin dynamics, which samples from a distribution by noisy gradient descent and is behind diffusion models.
\end{itemize}
We use the linear algebra of Lectures~1 and~2, the vector calculus of Lecture~3 (gradients, vector fields, the divergence and the divergence theorem), and the probability of Lecture~4 (Gaussian distributions and Markov chains).

\paragraph{Notation.} As in Lectures~3 and~4, vectors are bold, random variables are denoted by capital letters, $\mathbb{E}$ is the expectation, and $\log$ is the natural logarithm.

\section{Ordinary Differential Equations}\label{L5-sec:ode}

A differential equation describes how a state changes continuously in time: at every moment, the rate of change of the state is a given function of the state. This function is a vector field, and the state moves along its arrows, with the velocity given by the arrow at its current position. Differential equations appear in machine learning in the analysis of training dynamics, in neural networks with continuous depth, and, with randomness added, in diffusion models and sampling algorithms.

\begin{definition}[Ordinary differential equation]\label{L5-def:ode}
    Let $F: \mathbb{R}^n \to \mathbb{R}^n$ be a vector field (Lecture~3, Definition~\ref{L3-def:vector-field}). A solution of the ordinary differential equation (ODE) $\mathbf{x}' = F(\mathbf{x})$ with the initial value $\mathbf{x}(0) = \mathbf{x}_0$ is a differentiable curve $\mathbf{x}(t)$ with $\mathbf{x}(0) = \mathbf{x}_0$ and $\mathbf{x}'(t) = F(\mathbf{x}(t))$ for all $t$.
\end{definition}

Geometrically, a solution is a curve that is tangent to the vector field at each of its points: at the point $\mathbf{x}$, it moves with the velocity $F(\mathbf{x})$.

\begin{theorem}[Existence and uniqueness]\label{L5-thm:picard}
    If $F$ is Lipschitz continuous, i.e., $\norm{F(\mathbf{x}) - F(\mathbf{y})} \leq L\norm{\mathbf{x} - \mathbf{y}}$ for all $\mathbf{x}, \mathbf{y}$, then for every $\mathbf{x}_0$, the ODE $\mathbf{x}' = F(\mathbf{x})$ has exactly one solution with $\mathbf{x}(0) = \mathbf{x}_0$, defined for all $t \in \mathbb{R}$.
\end{theorem}

We do not prove this theorem (the Picard--Lindel\"of theorem). The idea is that $\mathbf{x}$ solves the problem if and only if $\mathbf{x}(t) = \mathbf{x}_0 + \int_0^t F(\mathbf{x}(s))\, ds$ for all $t$, so a solution is a fixed point of the map that sends a curve $\mathbf{x}$ to the right-hand side; on short time intervals, this map is a contraction, and the Banach fixed-point theorem of Appendix~\ref{A-chapter} applies.

\begin{example}\label{L5-ex:scalar-ode}
    \begin{enumerate}
        \item The solution of $x' = ax$ with $x(0) = x_0$ is $x(t) = e^{at}x_0$, since $\frac{d}{dt}e^{at}x_0 = ae^{at}x_0$. It grows exponentially if $a > 0$ and decays to $0$ if $a < 0$.
        \item The logistic equation $x' = x(1 - x)$ models growth with a limited capacity. By Example~\ref{L3-ex:sigmoid}, the sigmoid function satisfies $\sigma' = \sigma(1 - \sigma)$, so $x(t) = \sigma(t + c)$ is a solution, with $c$ chosen such that $\sigma(c) = x_0$. For $0 < x_0 < 1$, the solution increases and tends to $1$.
    \end{enumerate}
\end{example}

\section{Linear Systems}\label{L5-sec:linear-ode}


For a matrix $A \in \mathbb{R}^{n \times n}$, the linear ODE $\mathbf{x}' = A\mathbf{x}$ couples the coordinates, and eigenvectors decouple them (Lecture~1).

\begin{proposition}\label{L5-prop:linear-ode}
    Let $A = V\Lambda V^{-1}$ be diagonalizable, with eigenvalues $\lambda_1, \ldots, \lambda_n$ and eigenvectors $\mathbf{v}_1, \ldots, \mathbf{v}_n$ (the columns of $V$), and write $\mathbf{x}_0 = c_1\mathbf{v}_1 + \cdots + c_n\mathbf{v}_n$. Then the solution of $\mathbf{x}' = A\mathbf{x}$ with $\mathbf{x}(0) = \mathbf{x}_0$ is
    \[\mathbf{x}(t) = c_1e^{\lambda_1 t}\mathbf{v}_1 + \cdots + c_ne^{\lambda_n t}\mathbf{v}_n.\]
\end{proposition}

\begin{proof}
    Each term satisfies $\frac{d}{dt}c_ie^{\lambda_i t}\mathbf{v}_i = \lambda_ic_ie^{\lambda_i t}\mathbf{v}_i = A(c_ie^{\lambda_i t}\mathbf{v}_i)$, so the sum is a solution, and it has the right initial value. It is the only one by Theorem~\ref{L5-thm:picard}, since $\norm{A\mathbf{x} - A\mathbf{y}} \leq \norm{A}_2\norm{\mathbf{x} - \mathbf{y}}$. (Equivalently, the coordinates $\mathbf{y} = V^{-1}\mathbf{x}$ satisfy the decoupled equations $y_i' = \lambda_i y_i$.)
\end{proof}

More generally, the solution is $\mathbf{x}(t) = e^{At}\mathbf{x}_0$ with the matrix exponential $e^{At} = \sum_{k=0}^\infty \frac{(At)^k}{k!}$, which converges for every matrix and satisfies $\frac{d}{dt}e^{At} = Ae^{At}$; for diagonalizable $A$, $e^{At} = Ve^{\Lambda t}V^{-1}$ with $e^{\Lambda t} = \diag(e^{\lambda_1 t}, \ldots, e^{\lambda_n t})$. So the eigenvalues determine the long-term behavior: if all (complex) eigenvalues have negative real parts, every solution tends to $\mathbf{0}$, and if some eigenvalue has a positive real part, solutions starting in the direction of its eigenvector grow exponentially. This is the continuous-time analogue of the powers $P^t$ of a matrix, which govern gradient descent (Theorem~\ref{L3-thm:gd-quadratic}) and Markov chains (Lecture~4).

\begin{example}\label{L5-ex:linear-ode}
    \begin{enumerate}
        \item Let $A = \begin{pmatrix} -2 & 1 \\ 1 & -2 \end{pmatrix}$, with the eigenvalues $-1$ and $-3$ and the eigenvectors $\mathbf{v}_1 = (1, 1)^\top$ and $\mathbf{v}_2 = (1, -1)^\top$. For $\mathbf{x}_0 = (2, 0)^\top = \mathbf{v}_1 + \mathbf{v}_2$, the solution is $\mathbf{x}(t) = e^{-t}(1, 1)^\top + e^{-3t}(1, -1)^\top$; for example, $\mathbf{x}(1) \approx (0.418, 0.318)^\top$. The component along $\mathbf{v}_2$ decays faster, so all solutions approach the origin along the slow direction $\mathbf{v}_1$ (Figure~\ref{L5-fig:ode}).
        \item Let $A = \begin{pmatrix} 0 & -1 \\ 1 & 0 \end{pmatrix}$, with the complex eigenvalues $\pm i$. Then $e^{At} = \begin{pmatrix} \cos t & -\sin t \\ \sin t & \cos t \end{pmatrix}$ is the rotation by the angle $t$ (Lecture~2): one checks that this matrix is $I$ at $t = 0$ and that its derivative is $A$ times itself. The solutions move on circles. Indeed, for every skew-symmetric $A$, $\frac{d}{dt}\norm{\mathbf{x}}^2 = 2\mathbf{x}^\top A\mathbf{x} = 0$.
    \end{enumerate}
\end{example}

\begin{figure}[ht]
    \centering
    \begin{tikzpicture}[scale=1.15, >=Stealth]
        \draw[dashed, gray] (-2.3,-2.3) -- (2.3,2.3) node[above right, black] {$\mathbf{v}_1$};
        \draw[dashed, gray] (-2.3,2.3) -- (2.3,-2.3) node[below right, black] {$\mathbf{v}_2$};
        \foreach \x in {-2,-1.5,...,2}
            \foreach \y in {-2,-1.5,...,2}
                \draw[->, gray!70] (\x,\y) -- ++({0.09*(-2*\x + \y)},{0.09*(\x - 2*\y)});
        \foreach \a in {15,45,...,345}
            \draw[->, thick, blue!70!black] plot[domain=0:4, samples=60, variable=\t]
                ({(cos(\a) + sin(\a))*exp(-\t) + (cos(\a) - sin(\a))*exp(-3*\t)}, {(cos(\a) + sin(\a))*exp(-\t) - (cos(\a) - sin(\a))*exp(-3*\t)});
        \draw[->, gray] (-2.5,0) -- (2.6,0) node[right] {$x_1$};
        \draw[->, gray] (0,-2.5) -- (0,2.6) node[above] {$x_2$};
    \end{tikzpicture}
    \caption{The vector field $A\mathbf{x}$ (gray) and solutions of $\mathbf{x}' = A\mathbf{x}$ (blue) for $A = \left(\begin{smallmatrix} -2 & 1 \\ 1 & -2 \end{smallmatrix}\right)$, starting on the circle of radius $2$. The component along the eigenvector $\mathbf{v}_2$ (eigenvalue $-3$) dies out quickly, and the solutions approach the origin along $\mathbf{v}_1$ (eigenvalue $-1$).}
    \label{L5-fig:ode}
\end{figure}

\section{Gradient Flow and Euler's Method}\label{L5-sec:euler}


\begin{definition}[Gradient flow]\label{L5-def:gradient-flow}
    The gradient flow of a differentiable function $f: \mathbb{R}^n \to \mathbb{R}$ is the ODE $\mathbf{x}' = -\nabla f(\mathbf{x})$.
\end{definition}

\begin{proposition}\label{L5-prop:gradient-flow}
    Along every solution of the gradient flow, $\frac{d}{dt}f(\mathbf{x}(t)) = -\norm{\nabla f(\mathbf{x}(t))}^2 \leq 0$: the value of $f$ decreases, and strictly unless $\mathbf{x}(t)$ is a critical point.
\end{proposition}

\begin{proof}
    By the chain rule (Corollary~\ref{L3-cor:chain-special}), $\frac{d}{dt}f(\mathbf{x}(t)) = \nabla f(\mathbf{x}(t))^\top\mathbf{x}'(t) = -\norm{\nabla f(\mathbf{x}(t))}^2$.
\end{proof}

For the quadratic function $f(\mathbf{x}) = \frac{1}{2}\mathbf{x}^\top A\mathbf{x} - \mathbf{b}^\top\mathbf{x}$ with symmetric positive definite $A$, the error $\mathbf{e} = \mathbf{x} - \mathbf{x}^*$ satisfies $\mathbf{e}' = -A\mathbf{e}$, so by Proposition~\ref{L5-prop:linear-ode}, $\mathbf{e}(t) = \sum_i c_ie^{-\lambda_i t}\mathbf{q}_i$ in an orthonormal eigenbasis of $A$. The gradient flow converges to the minimum for every positive definite $A$, and the slowest direction decays like $e^{-\lambda_{\min}t}$.

To solve an ODE $\mathbf{x}' = F(\mathbf{x})$ numerically, Euler's method follows the tangent for a short time $h$:
\[\mathbf{x}_{k+1} = \mathbf{x}_k + hF(\mathbf{x}_k), \qquad \text{so that } \mathbf{x}_k \approx \mathbf{x}(kh).\]
For a smooth $F$, its error on a fixed time interval is proportional to $h$. For example, for $x' = -x$ with $x(0) = 1$, Euler's method gives $x_k = (1 - h)^k$, so it approximates $x(1) = e^{-1} \approx 0.368$ by $0.25$ for $h = 0.5$, and by $0.9^{10} \approx 0.349$ for $h = 0.1$.

For the gradient flow, Euler's method is $\mathbf{x}_{k+1} = \mathbf{x}_k - h\nabla f(\mathbf{x}_k)$: gradient descent is Euler's method for the gradient flow, with the step size $\eta = h$. This explains the results of Section~\ref{L3-sec:gd}. For the quadratic function, the flow converges for every positive definite $A$, but Euler's method multiplies the error in the direction $\mathbf{q}_i$ by $1 - h\lambda_i$, which is stable only if $h < \frac{2}{\lambda_{\max}}$. The directions of large curvature force small steps, while the directions of small curvature need the time $\frac{1}{\lambda_{\min}}$, so about $\kappa = \frac{\lambda_{\max}}{\lambda_{\min}}$ steps are needed. Such ODEs are called stiff, and the zig-zagging of Figure~\ref{L3-fig:gd} is an instability of the discretization, not of the flow.

\begin{remark}[ODEs in machine learning]
    \begin{enumerate}
        \item The layers of a residual network compute $\mathbf{x}_{l+1} = \mathbf{x}_l + f_l(\mathbf{x}_l)$, which is an Euler step of an ODE $\mathbf{x}' = f(\mathbf{x}, t)$ with $h = 1$. Neural ODEs replace the layers by an ODE solver, and compute gradients by solving an adjoint ODE backwards in time, which is reverse-mode differentiation (Section~\ref{L3-sec:autodiff}) in continuous time.
        \item Gradient descent with momentum is a discretization of the second-order ODE $\mathbf{x}'' + \gamma\mathbf{x}' + \nabla f(\mathbf{x}) = \mathbf{0}$ of a heavy ball with friction rolling in the landscape $f$; with $\mathbf{v} = \mathbf{x}'$, it is a first-order system for $(\mathbf{x}, \mathbf{v})$.
        \item Adding random noise to an ODE gives a stochastic differential equation. Langevin dynamics and diffusion models, which sample from probability distributions in this way, are discussed in Sections~\ref{L5-sec:sde} and~\ref{L5-sec:langevin}.
    \end{enumerate}
\end{remark}

\section{Flows, Volumes and Densities}\label{L5-sec:flows}

Instead of following a single solution, we can move all points at once. Let $\varphi_t(\mathbf{x}_0)$ be the solution of $\mathbf{x}' = F(\mathbf{x})$ with $\mathbf{x}(0) = \mathbf{x}_0$ at the time $t$. The map $\varphi_t$, the flow of $F$, moves every point along the vector field for the time $t$; for example, the flow of $F(\mathbf{x}) = A\mathbf{x}$ is $\varphi_t(\mathbf{x}) = e^{At}\mathbf{x}$. In Lecture~3, we saw that the divergence $\Div F = \tr DF$ is the net outflow of $F$ per unit volume. For the flow, this means that the divergence is the rate at which the flow expands volumes. For linear vector fields, the effect of the flow on volumes can be computed exactly.

\begin{proposition}[Liouville's formula]\label{L5-prop:liouville}
    For every diagonalizable $A \in \mathbb{R}^{n \times n}$, $\det e^{At} = e^{t\tr A}$. So the flow $\varphi_t(\mathbf{x}) = e^{At}\mathbf{x}$ of $F(\mathbf{x}) = A\mathbf{x}$ multiplies all volumes by $e^{t\Div F}$.
\end{proposition}

\begin{proof}
    With $A = V\Lambda V^{-1}$, $\det e^{At} = \det(V)\det(e^{\Lambda t})\det(V)^{-1} = e^{\lambda_1 t}\cdots e^{\lambda_n t} = e^{t(\lambda_1 + \cdots + \lambda_n)} = e^{t\tr A}$, and the linear map $e^{At}$ multiplies volumes by $|\det e^{At}|$ (Lecture~1).
\end{proof}

(The formula holds for every square matrix; Exercise~\ref{L5-ex5:nilpotent} checks it for a matrix that is not diagonalizable.) For a general vector field and a small $t$, the flow $\varphi_t(\mathbf{y}) \approx \mathbf{y} + tF(\mathbf{y})$ has approximately the Jacobian $I + tDF(\mathbf{y})$, whose determinant is $1 + t\Div F(\mathbf{y})$ up to terms of order $t^2$ (Lemma~\ref{L3-lem:det-identity}). So a small region near $\mathbf{x}$ with the volume $v$ has, after the short time $t$, approximately the volume $(1 + t\Div F(\mathbf{x}))v$:
\begin{center}
    \emph{$\Div F(\mathbf{x})$ is the rate at which the flow of $F$ expands volumes near $\mathbf{x}$.}
\end{center}
In particular, the flow of a vector field with $\Div F = 0$ preserves volumes.

\begin{example}\label{L5-ex:flow-volumes}
    \begin{enumerate}
        \item $F(\mathbf{x}) = \mathbf{x}$ has $\Div F = n$. Its flow $\varphi_t(\mathbf{x}) = e^t\mathbf{x}$ multiplies lengths by $e^t$ and volumes by $e^{nt}$.
        \item The rotation field $(-y, x)^\top$ has $\Div F = 0$, and its flow, the rotation by the angle $t$ (Example~\ref{L5-ex:linear-ode}), preserves areas. It is the continuous-time version of gradient descent--ascent for the game $\min_x\max_y xy$ (Lecture~3, Remark~\ref{L3-rem:vector-fields-ml}): the flow circles around the equilibrium forever, and Euler's method, i.e., gradient descent--ascent with a constant step size, spirals outwards (Exercise~\ref{L5-ex5:euler-skew}).
        \item For the field $A\mathbf{x}$ of Figure~\ref{L5-fig:ode}, $\Div F = \tr A = -4$, so the flow shrinks areas by the factor $e^{-4t}$, in agreement with the eigenvalues $-1$ and $-3$.
        \item The gradient flow $F = -\nabla f$ has $\Div F = -\Delta f = -\tr\nabla^2 f$. Near a strict local minimum, where the Hessian is positive definite, the gradient flow compresses volumes, and nearby solutions are drawn together.
    \end{enumerate}
\end{example}

\subsection{Transport of densities}

If the initial value of an ODE is random, then so is the solution at every time $t$, and the divergence also tells us how its probability density changes.

\begin{proposition}[Continuity equation]\label{L5-prop:continuity}
    Let the initial value of the ODE $\mathbf{x}' = F(\mathbf{x})$ be random, with the density $p_0$, and let $p_t$ be the density of $\mathbf{x}(t)$. Then, if all functions are smooth enough,
    \[\frac{\partial p_t}{\partial t} = -\Div(p_tF), \qquad \text{and along every solution,} \qquad \frac{d}{dt}\log p_t(\mathbf{x}(t)) = -\Div F(\mathbf{x}(t)).\]
\end{proposition}

\begin{proof}[Sketch of proof]
    Probability is neither created nor destroyed; it moves with the velocity $F$. So for a fixed region $V$, the probability $\int_V p_t\, d\mathbf{x}$ of $V$ changes by the flux of $p_tF$ out of $V$, and by the divergence theorem (Lecture~3, Theorem~\ref{L3-thm:divergence}),
    \[\frac{d}{dt}\int_V p_t\, d\mathbf{x} = -\int_{\partial V} p_tF^\top\mathbf{n}\, dS = -\int_V \Div(p_tF)\, d\mathbf{x}.\]
    Differentiating under the integral sign on the left, the integrals of $\frac{\partial p_t}{\partial t}$ and of $-\Div(p_tF)$ agree on every region $V$, so the continuous functions agree. For the second formula, the chain rule and $\Div(p_tF) = \nabla p_t^\top F + p_t\Div F$ give
    \begin{align*}
        \frac{d}{dt}p_t(\mathbf{x}(t)) &= \frac{\partial p_t}{\partial t}(\mathbf{x}(t)) + \nabla p_t(\mathbf{x}(t))^\top F(\mathbf{x}(t)) = -p_t(\mathbf{x}(t))\Div F(\mathbf{x}(t)),
    \end{align*}
    and we divide by $p_t(\mathbf{x}(t))$.
\end{proof}

The second formula is the continuous-time version of the change of variables formula of Remark~\ref{L3-rem:change-of-variables}: where the flow expands volumes ($\Div F > 0$), the density along the solution decreases, and conversely. Continuous normalizing flows, which are generative models, transform a Gaussian distribution by the flow of an ODE $\mathbf{x}' = F(\mathbf{x})$ whose vector field is a neural network, and compute densities by
\[\log p_T(\mathbf{x}(T)) = \log p_0(\mathbf{x}(0)) - \int_0^T \Div F(\mathbf{x}(t))\, dt,\]
which replaces the determinant of a Jacobian by the integral of a trace, estimated with Hutchinson's estimator (Lecture~3). If noise is added to the ODE, the continuity equation becomes the Fokker--Planck equation of Section~\ref{L5-sec:langevin}, which explains why Langevin dynamics samples from a given distribution.

\section{Stochastic Differential Equations}\label{L5-sec:sde}

Markov chains (Lecture~4) move in discrete steps. Many random processes change continuously in time: the position of a particle in a fluid, the price of a stock, the parameters of a network during training with stochastic gradients, or the gradually noised images of a diffusion model. Stochastic differential equations (SDEs) describe such processes by adding random noise to ordinary differential equations. The source of the noise is Brownian motion.

\begin{definition}[Brownian motion]\label{L5-def:brownian}
    A Brownian motion (or Wiener process) is a family of random variables $B(t)$, $t \geq 0$, with $B(0) = 0$, such that
    \begin{enumerate}
        \item for $0 \leq t_0 < t_1 < \cdots < t_k$, the increments $B(t_1) - B(t_0), \ldots, B(t_k) - B(t_{k-1})$ are independent;
        \item $B(t) - B(s) \sim \mathcal{N}(0, t - s)$ for $s < t$;
        \item the paths $t \mapsto B(t)$ are continuous.
    \end{enumerate}
\end{definition}

That such a process exists is a theorem of Norbert Wiener, which we do not prove. On the grid $t_k = kh$, Brownian motion is simulated exactly by the random walk $B(t_{k+1}) = B(t_k) + \sqrt{h}\,Z_k$ with i.i.d.\ $Z_k \sim \mathcal{N}(0, 1)$ (Figure~\ref{L5-fig:sde}). Conversely, by the central limit theorem, every random walk with many small independent steps looks like a Brownian motion from far away.

\begin{proposition}\label{L5-prop:brownian}
    $\E[B(t)] = 0$, $\Var(B(t)) = t$, and $\Cov(B(s), B(t)) = \min(s, t)$.
\end{proposition}

\begin{proof}
    The first two statements are property (2) with $s = 0$. For $s \leq t$, $\Cov(B(s), B(t)) = \Cov\big(B(s), B(s) + (B(t) - B(s))\big) = \Var(B(s)) + 0 = s$, since the increment $B(t) - B(s)$ is independent of $B(s) = B(s) - B(0)$.
\end{proof}

So Brownian motion is a Gaussian process with mean $0$ and covariance function $\min(s, t)$ (Section~\ref{L4-sec:gp}). Its typical displacement in a time $h$ is $\sqrt{h}$, which is much larger than $h$ for small $h$: the difference quotient $\frac{B(t + h) - B(t)}{h}$ has the variance $\frac{1}{h} \to \infty$. In fact, the paths of Brownian motion are continuous but nowhere differentiable, and ordinary calculus does not apply to them.

\begin{definition}[Stochastic differential equation]\label{L5-def:sde}
    An SDE $dX = \mu(X)\,dt + \sigma(X)\,dB$, with drift $\mu$ and diffusion coefficient $\sigma$, describes a process with $X(t) = X(0) + \int_0^t \mu(X(s))\, ds + \int_0^t \sigma(X(s))\, dB(s)$, where the last integral is a stochastic (It\^o) integral. Its Euler--Maruyama approximation with step $h$ is
    \[X_{k+1} = X_k + \mu(X_k)\,h + \sigma(X_k)\sqrt{h}\,Z_k, \qquad Z_k \sim \mathcal{N}(0, 1) \text{ i.i.d.}\]
\end{definition}

We do not define the stochastic integral; for our purposes, the Euler--Maruyama scheme, which converges to the solution as $h \to 0$, can serve as the definition. It is Euler's method for the ODE $x' = \mu(x)$ (Section~\ref{L5-sec:euler}), plus Gaussian noise whose standard deviation is proportional to $\sqrt{h}$, and it is a Markov chain with continuous states. For $\sigma = 0$, the SDE is an ODE.

\begin{example}[The Ornstein--Uhlenbeck process]\label{L5-ex:ou}
    The SDE $dX = -\theta X\,dt + \sigma\,dB$ with $\theta > 0$ pulls $X$ towards $0$ while the noise pushes it around. Its Euler--Maruyama approximation is $X_{k+1} = (1 - \theta h)X_k + \sigma\sqrt{h}\,Z_k$. Its mean satisfies $\E[X_{k+1}] = (1 - \theta h)\E[X_k]$, so it follows Euler's method for the ODE $x' = -\theta x$, with the solution $e^{-\theta t}X(0)$. If $X_k \sim \mathcal{N}(0, v)$ is independent of $Z_k$, then $X_{k+1} \sim \mathcal{N}\big(0, (1 - \theta h)^2v + \sigma^2h\big)$ by Corollary~\ref{L4-cor:gaussian-linear}. So for $0 < \theta h < 2$, the chain has the stationary distribution $\mathcal{N}(0, v_h)$ with
    \[v_h = \frac{\sigma^2h}{1 - (1 - \theta h)^2} = \frac{\sigma^2}{\theta(2 - \theta h)} \to \frac{\sigma^2}{2\theta} \qquad \text{as } h \to 0.\]
    For $\theta = 1$, $\sigma = \sqrt{2}$, the stationary distribution of the SDE is $\mathcal{N}(0, 1)$ (Figure~\ref{L5-fig:sde}); for $\theta = \sigma = 1$ and $h = 0.01$, $v_h \approx 0.5025$ instead of $0.5$.
\end{example}

\begin{figure}[ht]
    \centering
    \begin{tikzpicture}[x=1.2cm, y=0.48cm, >=Stealth]
        \draw[->, gray] (0,-6.2) -- (5.4,-6.2) node[right] {$t$};
        \draw[->, gray] (0,-6.2) -- (0,1.9) node[above] {$B(t)$};
        \draw[gray, dotted] (0,0) -- (5.2,0);
        \foreach \v in {-4,-2,0}
            \draw[gray] (0.06,\v) -- (-0.06,\v) node[left, font=\scriptsize] {$\v$};
        \foreach \v in {0,...,5}
            \draw[gray] (\v,-6.05) -- (\v,-6.35) node[below, font=\scriptsize] {$\v$};
        \draw[blue!70!black] plot coordinates {(0.00,0.000) (0.02,0.000) (0.04,0.042) (0.06,0.004) (0.08,-0.122) (0.10,-0.187) (0.12,-0.327) (0.14,-0.318) (0.16,-0.129) (0.18,-0.198) (0.20,-0.286) (0.22,-0.217) (0.24,-0.166) (0.26,-0.151) (0.28,-0.283) (0.30,-0.287) (0.32,-0.189) (0.34,-0.379) (0.36,-0.444) (0.38,-0.713) (0.40,-0.895) (0.42,-1.155) (0.44,-1.189) (0.46,-1.368) (0.48,-1.330) (0.50,-1.307) (0.52,-1.334) (0.54,-1.690) (0.56,-1.766) (0.58,-1.773) (0.60,-1.757) (0.62,-1.973) (0.64,-2.041) (0.66,-2.179) (0.68,-2.293) (0.70,-2.143) (0.72,-2.258) (0.74,-2.262) (0.76,-2.137) (0.78,-2.220) (0.80,-2.236) (0.82,-2.220) (0.84,-2.211) (0.86,-2.384) (0.88,-2.373) (0.90,-2.181) (0.92,-2.400) (0.94,-2.278) (0.96,-2.262) (0.98,-2.352) (1.00,-2.069) (1.02,-1.962) (1.04,-2.131) (1.06,-2.121) (1.08,-2.039) (1.10,-2.066) (1.12,-1.969) (1.14,-1.979) (1.16,-1.884) (1.18,-1.681) (1.20,-1.776) (1.22,-1.748) (1.24,-1.813) (1.26,-1.795) (1.28,-1.963) (1.30,-2.045) (1.32,-2.073) (1.34,-1.946) (1.36,-1.784) (1.38,-1.971) (1.40,-2.083) (1.42,-1.992) (1.44,-2.274) (1.46,-2.339) (1.48,-2.353) (1.50,-2.175) (1.52,-2.078) (1.54,-2.124) (1.56,-2.176) (1.58,-2.211) (1.60,-1.996) (1.62,-2.056) (1.64,-2.099) (1.66,-2.049) (1.68,-2.067) (1.70,-2.094) (1.72,-2.252) (1.74,-2.254) (1.76,-2.316) (1.78,-2.151) (1.80,-2.059) (1.82,-2.062) (1.84,-1.968) (1.86,-2.016) (1.88,-1.867) (1.90,-1.868) (1.92,-1.785) (1.94,-1.968) (1.96,-1.919) (1.98,-2.158) (2.00,-2.446) (2.02,-2.489) (2.04,-2.616) (2.06,-2.593) (2.08,-2.275) (2.10,-2.393) (2.12,-2.481) (2.14,-2.452) (2.16,-2.382) (2.18,-2.407) (2.20,-2.436) (2.22,-2.337) (2.24,-2.264) (2.26,-2.410) (2.28,-2.421) (2.30,-2.416) (2.32,-2.565) (2.34,-2.528) (2.36,-2.650) (2.38,-2.512) (2.40,-2.485) (2.42,-2.472) (2.44,-2.556) (2.46,-2.573) (2.48,-2.855) (2.50,-3.015) (2.52,-2.964) (2.54,-3.265) (2.56,-3.145) (2.58,-3.392) (2.60,-3.285) (2.62,-3.405) (2.64,-3.295) (2.66,-3.276) (2.68,-3.493) (2.70,-3.317) (2.72,-3.113) (2.74,-3.122) (2.76,-3.161) (2.78,-3.183) (2.80,-3.321) (2.82,-3.166) (2.84,-3.243) (2.86,-3.250) (2.88,-3.362) (2.90,-3.451) (2.92,-3.631) (2.94,-3.454) (2.96,-3.475) (2.98,-3.339) (3.00,-3.337) (3.02,-3.435) (3.04,-3.481) (3.06,-3.561) (3.08,-3.559) (3.10,-3.613) (3.12,-3.655) (3.14,-3.850) (3.16,-3.964) (3.18,-3.730) (3.20,-3.825) (3.22,-3.974) (3.24,-3.926) (3.26,-3.727) (3.28,-3.933) (3.30,-3.962) (3.32,-4.052) (3.34,-4.301) (3.36,-4.197) (3.38,-4.200) (3.40,-4.190) (3.42,-4.297) (3.44,-4.232) (3.46,-4.309) (3.48,-4.329) (3.50,-4.485) (3.52,-4.657) (3.54,-4.469) (3.56,-4.540) (3.58,-4.499) (3.60,-4.504) (3.62,-4.566) (3.64,-4.638) (3.66,-4.549) (3.68,-4.592) (3.70,-4.613) (3.72,-4.610) (3.74,-4.444) (3.76,-4.347) (3.78,-4.293) (3.80,-4.373) (3.82,-4.568) (3.84,-4.434) (3.86,-4.297) (3.88,-4.317) (3.90,-4.241) (3.92,-4.130) (3.94,-4.013) (3.96,-3.882) (3.98,-3.947) (4.00,-3.732) (4.02,-3.909) (4.04,-3.787) (4.06,-3.717) (4.08,-3.593) (4.10,-3.328) (4.12,-3.118) (4.14,-3.280) (4.16,-3.519) (4.18,-3.403) (4.20,-3.547) (4.22,-3.548) (4.24,-3.430) (4.26,-3.662) (4.28,-3.960) (4.30,-3.924) (4.32,-3.918) (4.34,-3.952) (4.36,-3.947) (4.38,-4.069) (4.40,-4.283) (4.42,-4.306) (4.44,-4.444) (4.46,-4.676) (4.48,-4.604) (4.50,-4.613) (4.52,-4.556) (4.54,-4.696) (4.56,-4.789) (4.58,-4.930) (4.60,-5.055) (4.62,-5.028) (4.64,-5.138) (4.66,-5.088) (4.68,-5.040) (4.70,-4.754) (4.72,-4.951) (4.74,-4.825) (4.76,-4.838) (4.78,-4.840) (4.80,-5.045) (4.82,-5.110) (4.84,-5.005) (4.86,-5.016) (4.88,-5.005) (4.90,-5.046) (4.92,-4.883) (4.94,-4.886) (4.96,-5.197) (4.98,-5.295) (5.00,-5.573)};
        \draw[red!70!black] plot coordinates {(0.00,0.000) (0.02,-0.460) (0.04,-0.535) (0.06,-0.346) (0.08,-0.340) (0.10,-0.505) (0.12,-0.638) (0.14,-0.479) (0.16,-0.456) (0.18,-0.449) (0.20,-0.457) (0.22,-0.452) (0.24,-0.338) (0.26,-0.260) (0.28,-0.229) (0.30,-0.376) (0.32,-0.304) (0.34,-0.401) (0.36,-0.246) (0.38,-0.426) (0.40,-0.445) (0.42,-0.447) (0.44,-0.634) (0.46,-0.390) (0.48,-0.184) (0.50,-0.249) (0.52,-0.140) (0.54,-0.087) (0.56,-0.456) (0.58,-0.421) (0.60,-0.430) (0.62,-0.418) (0.64,-0.570) (0.66,-0.608) (0.68,-0.633) (0.70,-0.465) (0.72,-0.418) (0.74,-0.419) (0.76,-0.203) (0.78,-0.281) (0.80,-0.336) (0.82,-0.593) (0.84,-0.371) (0.86,-0.235) (0.88,-0.105) (0.90,-0.011) (0.92,0.005) (0.94,0.035) (0.96,-0.000) (0.98,-0.029) (1.00,-0.021) (1.02,0.192) (1.04,0.271) (1.06,0.263) (1.08,0.181) (1.10,0.091) (1.12,0.318) (1.14,0.389) (1.16,0.399) (1.18,0.350) (1.20,0.193) (1.22,0.184) (1.24,0.307) (1.26,0.252) (1.28,0.220) (1.30,0.188) (1.32,0.204) (1.34,-0.021) (1.36,-0.055) (1.38,-0.176) (1.40,-0.050) (1.42,-0.159) (1.44,-0.078) (1.46,0.138) (1.48,0.093) (1.50,0.008) (1.52,0.035) (1.54,0.035) (1.56,-0.105) (1.58,-0.040) (1.60,0.245) (1.62,0.208) (1.64,0.180) (1.66,0.032) (1.68,0.077) (1.70,-0.099) (1.72,-0.256) (1.74,-0.075) (1.76,-0.203) (1.78,-0.050) (1.80,0.165) (1.82,0.202) (1.84,0.280) (1.86,0.557) (1.88,0.529) (1.90,0.445) (1.92,0.253) (1.94,0.259) (1.96,0.469) (1.98,0.604) (2.00,0.471) (2.02,0.350) (2.04,0.279) (2.06,0.320) (2.08,0.291) (2.10,0.321) (2.12,0.363) (2.14,0.321) (2.16,0.315) (2.18,0.345) (2.20,0.333) (2.22,0.404) (2.24,0.669) (2.26,0.752) (2.28,0.760) (2.30,0.522) (2.32,0.577) (2.34,0.301) (2.36,0.102) (2.38,0.223) (2.40,0.323) (2.42,0.302) (2.44,0.060) (2.46,0.007) (2.48,-0.089) (2.50,0.001) (2.52,0.321) (2.54,0.351) (2.56,0.241) (2.58,0.075) (2.60,0.068) (2.62,0.043) (2.64,-0.120) (2.66,-0.104) (2.68,-0.267) (2.70,-0.109) (2.72,0.041) (2.74,0.194) (2.76,0.127) (2.78,0.200) (2.80,0.181) (2.82,0.126) (2.84,0.078) (2.86,-0.105) (2.88,-0.310) (2.90,-0.197) (2.92,-0.224) (2.94,-0.194) (2.96,-0.052) (2.98,-0.297) (3.00,-0.408) (3.02,-0.383) (3.04,-0.328) (3.06,-0.381) (3.08,-0.236) (3.10,-0.206) (3.12,-0.377) (3.14,-0.509) (3.16,-0.395) (3.18,-0.330) (3.20,-0.598) (3.22,-0.407) (3.24,-0.323) (3.26,-0.133) (3.28,-0.187) (3.30,-0.229) (3.32,-0.388) (3.34,-0.030) (3.36,-0.054) (3.38,0.170) (3.40,0.079) (3.42,0.102) (3.44,-0.135) (3.46,-0.189) (3.48,-0.050) (3.50,-0.227) (3.52,-0.075) (3.54,-0.027) (3.56,-0.175) (3.58,-0.246) (3.60,-0.311) (3.62,-0.318) (3.64,-0.394) (3.66,-0.511) (3.68,-0.554) (3.70,-0.699) (3.72,-0.881) (3.74,-0.888) (3.76,-0.763) (3.78,-0.980) (3.80,-0.979) (3.82,-1.071) (3.84,-1.209) (3.86,-1.088) (3.88,-1.162) (3.90,-0.950) (3.92,-1.060) (3.94,-1.005) (3.96,-1.038) (3.98,-1.144) (4.00,-1.061) (4.02,-1.083) (4.04,-0.998) (4.06,-1.004) (4.08,-1.158) (4.10,-1.172) (4.12,-1.165) (4.14,-1.030) (4.16,-1.158) (4.18,-1.163) (4.20,-1.407) (4.22,-1.315) (4.24,-1.468) (4.26,-1.723) (4.28,-1.731) (4.30,-1.575) (4.32,-1.791) (4.34,-1.944) (4.36,-2.050) (4.38,-2.209) (4.40,-2.156) (4.42,-2.270) (4.44,-2.372) (4.46,-2.289) (4.48,-2.396) (4.50,-2.335) (4.52,-2.472) (4.54,-2.644) (4.56,-2.903) (4.58,-2.640) (4.60,-2.685) (4.62,-2.651) (4.64,-2.655) (4.66,-2.633) (4.68,-2.626) (4.70,-2.356) (4.72,-2.503) (4.74,-2.723) (4.76,-2.866) (4.78,-3.055) (4.80,-2.949) (4.82,-2.833) (4.84,-2.969) (4.86,-3.166) (4.88,-3.216) (4.90,-3.019) (4.92,-3.418) (4.94,-3.344) (4.96,-3.496) (4.98,-3.349) (5.00,-3.501)};
        \draw[orange!85!black] plot coordinates {(0.00,0.000) (0.02,-0.040) (0.04,-0.253) (0.06,-0.392) (0.08,-0.196) (0.10,-0.080) (0.12,-0.136) (0.14,-0.259) (0.16,-0.527) (0.18,-0.583) (0.20,-0.587) (0.22,-0.599) (0.24,-0.612) (0.26,-0.771) (0.28,-0.780) (0.30,-0.786) (0.32,-0.603) (0.34,-0.339) (0.36,-0.359) (0.38,-0.467) (0.40,-0.476) (0.42,-0.562) (0.44,-0.667) (0.46,-0.676) (0.48,-0.823) (0.50,-0.737) (0.52,-0.752) (0.54,-0.717) (0.56,-0.743) (0.58,-0.845) (0.60,-0.980) (0.62,-1.013) (0.64,-1.091) (0.66,-1.058) (0.68,-1.058) (0.70,-1.251) (0.72,-1.241) (0.74,-1.431) (0.76,-1.518) (0.78,-1.560) (0.80,-1.854) (0.82,-1.841) (0.84,-1.819) (0.86,-1.842) (0.88,-1.902) (0.90,-1.955) (0.92,-2.093) (0.94,-2.131) (0.96,-2.209) (0.98,-2.196) (1.00,-2.366) (1.02,-2.333) (1.04,-2.313) (1.06,-2.333) (1.08,-2.395) (1.10,-2.317) (1.12,-2.552) (1.14,-2.487) (1.16,-2.453) (1.18,-2.413) (1.20,-2.358) (1.22,-2.451) (1.24,-2.487) (1.26,-2.397) (1.28,-2.336) (1.30,-2.307) (1.32,-2.521) (1.34,-2.445) (1.36,-2.280) (1.38,-2.137) (1.40,-2.104) (1.42,-2.324) (1.44,-2.191) (1.46,-2.212) (1.48,-2.570) (1.50,-2.517) (1.52,-2.728) (1.54,-2.911) (1.56,-3.001) (1.58,-2.821) (1.60,-2.873) (1.62,-2.835) (1.64,-2.588) (1.66,-2.362) (1.68,-2.377) (1.70,-2.411) (1.72,-2.589) (1.74,-2.688) (1.76,-2.628) (1.78,-2.572) (1.80,-2.556) (1.82,-2.415) (1.84,-2.524) (1.86,-2.532) (1.88,-2.429) (1.90,-2.346) (1.92,-2.195) (1.94,-2.139) (1.96,-2.183) (1.98,-2.131) (2.00,-2.273) (2.02,-2.505) (2.04,-2.423) (2.06,-2.431) (2.08,-2.387) (2.10,-2.627) (2.12,-2.679) (2.14,-2.764) (2.16,-2.886) (2.18,-3.205) (2.20,-3.252) (2.22,-3.125) (2.24,-3.071) (2.26,-3.156) (2.28,-3.158) (2.30,-3.051) (2.32,-3.442) (2.34,-3.459) (2.36,-3.383) (2.38,-3.286) (2.40,-3.046) (2.42,-2.885) (2.44,-2.841) (2.46,-2.798) (2.48,-2.687) (2.50,-2.763) (2.52,-2.769) (2.54,-2.641) (2.56,-2.364) (2.58,-2.386) (2.60,-2.393) (2.62,-2.365) (2.64,-2.175) (2.66,-2.179) (2.68,-1.972) (2.70,-2.108) (2.72,-2.135) (2.74,-2.163) (2.76,-2.051) (2.78,-1.904) (2.80,-2.117) (2.82,-2.247) (2.84,-2.199) (2.86,-2.292) (2.88,-2.507) (2.90,-2.360) (2.92,-2.291) (2.94,-2.221) (2.96,-2.288) (2.98,-2.143) (3.00,-2.177) (3.02,-2.021) (3.04,-2.150) (3.06,-2.271) (3.08,-2.242) (3.10,-2.342) (3.12,-2.246) (3.14,-2.209) (3.16,-2.340) (3.18,-2.329) (3.20,-2.379) (3.22,-2.250) (3.24,-2.339) (3.26,-2.401) (3.28,-2.230) (3.30,-1.913) (3.32,-1.631) (3.34,-1.622) (3.36,-1.591) (3.38,-1.374) (3.40,-1.391) (3.42,-1.530) (3.44,-1.513) (3.46,-1.449) (3.48,-1.566) (3.50,-1.799) (3.52,-2.002) (3.54,-1.908) (3.56,-2.016) (3.58,-2.036) (3.60,-2.005) (3.62,-1.918) (3.64,-1.965) (3.66,-1.895) (3.68,-2.021) (3.70,-2.072) (3.72,-2.217) (3.74,-2.057) (3.76,-2.060) (3.78,-2.165) (3.80,-2.215) (3.82,-2.246) (3.84,-2.147) (3.86,-2.373) (3.88,-2.519) (3.90,-2.573) (3.92,-2.215) (3.94,-2.080) (3.96,-2.095) (3.98,-1.995) (4.00,-1.704) (4.02,-1.737) (4.04,-1.789) (4.06,-1.617) (4.08,-1.547) (4.10,-1.452) (4.12,-1.524) (4.14,-1.252) (4.16,-1.010) (4.18,-0.930) (4.20,-0.834) (4.22,-1.120) (4.24,-1.030) (4.26,-1.058) (4.28,-0.996) (4.30,-0.900) (4.32,-0.948) (4.34,-1.187) (4.36,-1.135) (4.38,-1.240) (4.40,-1.287) (4.42,-1.372) (4.44,-1.420) (4.46,-1.747) (4.48,-1.575) (4.50,-1.539) (4.52,-1.382) (4.54,-1.102) (4.56,-1.099) (4.58,-1.354) (4.60,-1.480) (4.62,-1.651) (4.64,-1.722) (4.66,-1.710) (4.68,-1.993) (4.70,-1.945) (4.72,-2.159) (4.74,-2.117) (4.76,-2.132) (4.78,-2.176) (4.80,-2.187) (4.82,-2.263) (4.84,-2.350) (4.86,-2.588) (4.88,-2.592) (4.90,-2.331) (4.92,-2.051) (4.94,-1.864) (4.96,-1.764) (4.98,-1.860) (5.00,-1.656)};
    \end{tikzpicture}
    \quad
    \begin{tikzpicture}[x=1.2cm, y=0.62cm, >=Stealth]
        \draw[->, gray] (0,-3.2) -- (5.4,-3.2) node[right] {$t$};
        \draw[->, gray] (0,-3.2) -- (0,3.3) node[above] {$X(t)$};
        \foreach \v in {-2,0,2} {
            \draw[gray, dotted] (0,\v) -- (5.2,\v);
            \draw[gray] (0.06,\v) -- (-0.06,\v) node[left, font=\scriptsize] {$\v$};
        }
        \foreach \v in {0,...,5}
            \draw[gray] (\v,-3.08) -- (\v,-3.32) node[below, font=\scriptsize] {$\v$};
        \draw[blue!70!black, dashed, thin] plot[domain=0:5, samples=60] (\x, {2.5*exp(-\x)});
        \draw[red!70!black, dashed, thin] plot[domain=0:5, samples=60] (\x, {-2.0*exp(-\x)});
        \draw[orange!85!black, dashed, thin] plot[domain=0:5, samples=60] (\x, {0*exp(-\x)});
        \draw[blue!70!black] plot coordinates {(0.00,2.500) (0.02,2.439) (0.04,2.376) (0.06,2.270) (0.08,2.243) (0.10,2.112) (0.12,2.053) (0.14,1.795) (0.16,1.684) (0.18,2.106) (0.20,2.050) (0.22,1.962) (0.24,2.029) (0.26,2.130) (0.28,1.865) (0.30,1.786) (0.32,1.934) (0.34,1.949) (0.36,1.934) (0.38,2.207) (0.40,2.027) (0.42,2.003) (0.44,1.859) (0.46,2.120) (0.48,1.687) (0.50,1.519) (0.52,1.383) (0.54,1.488) (0.56,1.580) (0.58,1.827) (0.60,1.476) (0.62,1.596) (0.64,1.506) (0.66,1.342) (0.68,1.423) (0.70,1.211) (0.72,0.771) (0.74,0.682) (0.76,0.369) (0.78,0.232) (0.80,0.301) (0.82,0.358) (0.84,0.668) (0.86,0.615) (0.88,0.296) (0.90,0.138) (0.92,-0.048) (0.94,-0.291) (0.96,-0.198) (0.98,-0.323) (1.00,-0.712) (1.02,-0.558) (1.04,-0.570) (1.06,-0.487) (1.08,-0.456) (1.10,-0.320) (1.12,-0.306) (1.14,-0.053) (1.16,0.033) (1.18,0.111) (1.20,0.191) (1.22,-0.105) (1.24,-0.135) (1.26,-0.185) (1.28,-0.139) (1.30,-0.408) (1.32,-0.073) (1.34,-0.051) (1.36,-0.292) (1.38,-0.628) (1.40,-0.672) (1.42,-0.676) (1.44,-0.807) (1.46,-0.772) (1.48,-0.885) (1.50,-0.757) (1.52,-0.887) (1.54,-0.876) (1.56,-0.663) (1.58,-0.136) (1.60,-0.334) (1.62,-0.421) (1.64,-0.580) (1.66,-0.412) (1.68,-0.633) (1.70,-0.717) (1.72,-0.709) (1.74,-0.890) (1.76,-1.064) (1.78,-1.138) (1.80,-1.535) (1.82,-1.794) (1.84,-1.840) (1.86,-1.774) (1.88,-1.776) (1.90,-2.095) (1.92,-2.146) (1.94,-1.943) (1.96,-1.793) (1.98,-1.773) (2.00,-1.915) (2.02,-1.750) (2.04,-1.831) (2.06,-2.028) (2.08,-2.148) (2.10,-1.816) (2.12,-1.735) (2.14,-1.469) (2.16,-1.535) (2.18,-1.317) (2.20,-1.411) (2.22,-1.414) (2.24,-0.889) (2.26,-0.717) (2.28,-0.803) (2.30,-0.804) (2.32,-0.723) (2.34,-0.466) (2.36,-0.555) (2.38,-0.892) (2.40,-0.930) (2.42,-0.908) (2.44,-0.868) (2.46,-0.588) (2.48,-0.513) (2.50,-0.340) (2.52,-0.553) (2.54,-0.368) (2.56,0.058) (2.58,0.212) (2.60,0.258) (2.62,0.284) (2.64,0.635) (2.66,0.437) (2.68,0.406) (2.70,0.490) (2.72,0.629) (2.74,0.529) (2.76,0.580) (2.78,0.513) (2.80,0.527) (2.82,0.490) (2.84,0.252) (2.86,0.242) (2.88,0.413) (2.90,0.211) (2.92,0.159) (2.94,0.289) (2.96,0.069) (2.98,0.104) (3.00,-0.110) (3.02,0.119) (3.04,0.579) (3.06,0.972) (3.08,0.909) (3.10,1.039) (3.12,1.042) (3.14,1.042) (3.16,1.330) (3.18,1.040) (3.20,1.230) (3.22,1.196) (3.24,1.454) (3.26,1.462) (3.28,1.298) (3.30,1.328) (3.32,1.448) (3.34,1.427) (3.36,1.496) (3.38,1.361) (3.40,0.907) (3.42,1.069) (3.44,1.188) (3.46,1.194) (3.48,1.183) (3.50,1.367) (3.52,1.248) (3.54,1.082) (3.56,1.023) (3.58,1.240) (3.60,0.938) (3.62,1.157) (3.64,1.006) (3.66,0.766) (3.68,1.003) (3.70,0.963) (3.72,0.684) (3.74,0.599) (3.76,0.773) (3.78,0.996) (3.80,0.890) (3.82,0.954) (3.84,1.078) (3.86,0.927) (3.88,0.980) (3.90,0.954) (3.92,0.827) (3.94,0.712) (3.96,0.712) (3.98,0.703) (4.00,0.576) (4.02,0.479) (4.04,0.692) (4.06,0.721) (4.08,0.884) (4.10,1.107) (4.12,1.202) (4.14,1.632) (4.16,1.435) (4.18,1.568) (4.20,1.473) (4.22,1.814) (4.24,2.118) (4.26,1.685) (4.28,1.458) (4.30,1.561) (4.32,1.688) (4.34,1.802) (4.36,1.751) (4.38,1.807) (4.40,1.902) (4.42,1.848) (4.44,2.017) (4.46,1.524) (4.48,1.621) (4.50,1.382) (4.52,1.546) (4.54,1.469) (4.56,1.262) (4.58,1.311) (4.60,1.103) (4.62,0.898) (4.64,0.567) (4.66,0.550) (4.68,0.639) (4.70,0.830) (4.72,0.785) (4.74,0.979) (4.76,0.963) (4.78,0.925) (4.80,1.022) (4.82,1.213) (4.84,1.120) (4.86,1.049) (4.88,0.996) (4.90,0.993) (4.92,0.793) (4.94,0.982) (4.96,0.883) (4.98,0.958) (5.00,0.773)};
        \draw[red!70!black] plot coordinates {(0.00,-2.000) (0.02,-1.888) (0.04,-1.772) (0.06,-1.821) (0.08,-1.380) (0.10,-1.278) (0.12,-0.897) (0.14,-0.688) (0.16,-0.806) (0.18,-0.867) (0.20,-0.762) (0.22,-0.735) (0.24,-0.710) (0.26,-0.753) (0.28,-1.100) (0.30,-1.123) (0.32,-1.544) (0.34,-1.439) (0.36,-1.555) (0.38,-1.667) (0.40,-1.678) (0.42,-1.590) (0.44,-1.844) (0.46,-2.157) (0.48,-2.327) (0.50,-2.689) (0.52,-2.829) (0.54,-2.454) (0.56,-2.616) (0.58,-2.433) (0.60,-2.659) (0.62,-2.546) (0.64,-2.559) (0.66,-2.520) (0.68,-2.356) (0.70,-1.957) (0.72,-1.879) (0.74,-1.817) (0.76,-1.975) (0.78,-1.819) (0.80,-1.832) (0.82,-1.629) (0.84,-1.605) (0.86,-1.225) (0.88,-1.597) (0.90,-1.624) (0.92,-1.415) (0.94,-1.457) (0.96,-1.586) (0.98,-1.608) (1.00,-1.852) (1.02,-1.791) (1.04,-1.267) (1.06,-1.013) (1.08,-1.214) (1.10,-1.365) (1.12,-1.418) (1.14,-1.189) (1.16,-1.330) (1.18,-1.441) (1.20,-1.235) (1.22,-1.038) (1.24,-1.092) (1.26,-1.293) (1.28,-1.577) (1.30,-1.686) (1.32,-2.098) (1.34,-1.906) (1.36,-1.994) (1.38,-1.858) (1.40,-1.447) (1.42,-1.184) (1.44,-1.390) (1.46,-1.188) (1.48,-0.933) (1.50,-1.064) (1.52,-1.233) (1.54,-1.230) (1.56,-1.526) (1.58,-1.201) (1.60,-1.658) (1.62,-1.847) (1.64,-1.864) (1.66,-1.872) (1.68,-1.801) (1.70,-1.711) (1.72,-1.719) (1.74,-1.457) (1.76,-1.856) (1.78,-1.819) (1.80,-1.926) (1.82,-1.861) (1.84,-1.779) (1.86,-1.926) (1.88,-2.016) (1.90,-1.817) (1.92,-1.711) (1.94,-1.813) (1.96,-1.368) (1.98,-0.879) (2.00,-1.154) (2.02,-1.071) (2.04,-0.547) (2.06,-0.380) (2.08,-0.328) (2.10,-0.363) (2.12,-0.464) (2.14,-0.497) (2.16,-0.597) (2.18,-0.436) (2.20,-0.507) (2.22,-0.585) (2.24,-0.814) (2.26,-0.808) (2.28,-0.970) (2.30,-0.987) (2.32,-0.759) (2.34,-0.671) (2.36,-0.556) (2.38,-0.474) (2.40,-0.453) (2.42,-0.469) (2.44,-0.521) (2.46,-0.359) (2.48,-0.569) (2.50,-0.289) (2.52,-0.277) (2.54,-0.421) (2.56,-0.510) (2.58,-0.635) (2.60,-0.591) (2.62,-0.722) (2.64,-0.480) (2.66,-0.626) (2.68,-1.068) (2.70,-1.193) (2.72,-1.571) (2.74,-1.547) (2.76,-1.305) (2.78,-1.149) (2.80,-1.394) (2.82,-1.518) (2.84,-1.132) (2.86,-1.046) (2.88,-1.024) (2.90,-0.792) (2.92,-0.286) (2.94,-0.021) (2.96,0.007) (2.98,0.078) (3.00,0.200) (3.02,0.074) (3.04,-0.140) (3.06,-0.126) (3.08,-0.313) (3.10,-0.319) (3.12,-0.294) (3.14,0.180) (3.16,0.007) (3.18,-0.018) (3.20,-0.050) (3.22,0.038) (3.24,0.247) (3.26,0.144) (3.28,-0.024) (3.30,-0.351) (3.32,-0.529) (3.34,-0.412) (3.36,-0.394) (3.38,-0.589) (3.40,-0.652) (3.42,-0.633) (3.44,-0.519) (3.46,-0.628) (3.48,-0.665) (3.50,-1.016) (3.52,-1.226) (3.54,-0.878) (3.56,-1.298) (3.58,-1.339) (3.60,-1.269) (3.62,-1.318) (3.64,-1.453) (3.66,-1.437) (3.68,-1.409) (3.70,-1.397) (3.72,-1.745) (3.74,-1.739) (3.76,-1.875) (3.78,-1.948) (3.80,-1.863) (3.82,-1.698) (3.84,-1.484) (3.86,-1.554) (3.88,-1.338) (3.90,-1.077) (3.92,-0.828) (3.94,-0.533) (3.96,-0.552) (3.98,-0.576) (4.00,-0.399) (4.02,-0.664) (4.04,-0.609) (4.06,-0.703) (4.08,-0.763) (4.10,-1.096) (4.12,-1.252) (4.14,-1.231) (4.16,-1.029) (4.18,-0.810) (4.20,-0.810) (4.22,-0.832) (4.24,-0.981) (4.26,-0.881) (4.28,-0.913) (4.30,-0.774) (4.32,-0.408) (4.34,-0.406) (4.36,-0.697) (4.38,-0.856) (4.40,-1.130) (4.42,-1.347) (4.44,-1.059) (4.46,-0.992) (4.48,-1.277) (4.50,-1.120) (4.52,-0.846) (4.54,-0.901) (4.56,-1.018) (4.58,-1.065) (4.60,-0.987) (4.62,-0.839) (4.64,-0.586) (4.66,-0.332) (4.68,-0.088) (4.70,0.186) (4.72,0.315) (4.74,0.002) (4.76,-0.024) (4.78,0.039) (4.80,-0.038) (4.82,0.145) (4.84,0.070) (4.86,-0.117) (4.88,0.174) (4.90,0.304) (4.92,0.342) (4.94,0.520) (4.96,0.723) (4.98,0.776) (5.00,0.269)};
        \draw[orange!85!black] plot coordinates {(0.00,0.000) (0.02,-0.135) (0.04,-0.224) (0.06,-0.416) (0.08,-0.368) (0.10,-0.122) (0.12,-0.217) (0.14,-0.439) (0.16,-0.025) (0.18,-0.115) (0.20,-0.359) (0.22,-0.304) (0.24,-0.213) (0.26,-0.350) (0.28,-0.185) (0.30,-0.279) (0.32,-0.458) (0.34,-0.413) (0.36,-0.250) (0.38,-0.372) (0.40,-0.299) (0.42,-0.311) (0.44,0.260) (0.46,0.113) (0.48,0.388) (0.50,0.369) (0.52,0.337) (0.54,0.474) (0.56,0.643) (0.58,0.884) (0.60,0.932) (0.62,0.793) (0.64,0.670) (0.66,0.758) (0.68,0.859) (0.70,1.122) (0.72,1.183) (0.74,1.372) (0.76,1.647) (0.78,1.647) (0.80,1.317) (0.82,1.055) (0.84,0.747) (0.86,1.050) (0.88,0.859) (0.90,1.089) (0.92,1.184) (0.94,1.503) (0.96,1.674) (0.98,1.620) (1.00,1.548) (1.02,1.534) (1.04,1.538) (1.06,1.401) (1.08,1.367) (1.10,1.661) (1.12,1.287) (1.14,1.313) (1.16,1.106) (1.18,1.121) (1.20,1.267) (1.22,1.230) (1.24,1.360) (1.26,1.015) (1.28,1.216) (1.30,1.071) (1.32,1.174) (1.34,1.187) (1.36,0.646) (1.38,0.482) (1.40,0.516) (1.42,0.817) (1.44,0.858) (1.46,0.892) (1.48,0.592) (1.50,0.866) (1.52,1.210) (1.54,1.192) (1.56,1.124) (1.58,0.778) (1.60,0.939) (1.62,1.461) (1.64,1.575) (1.66,1.796) (1.68,1.869) (1.70,1.634) (1.72,1.869) (1.74,1.585) (1.76,1.511) (1.78,1.536) (1.80,1.682) (1.82,1.732) (1.84,1.774) (1.86,1.587) (1.88,1.297) (1.90,1.283) (1.92,1.116) (1.94,0.832) (1.96,0.562) (1.98,0.453) (2.00,0.073) (2.02,-0.198) (2.04,-0.521) (2.06,-0.474) (2.08,-0.383) (2.10,0.026) (2.12,-0.274) (2.14,-0.404) (2.16,-0.214) (2.18,-0.253) (2.20,-0.313) (2.22,0.035) (2.24,-0.033) (2.26,-0.264) (2.28,-0.522) (2.30,-0.445) (2.32,-0.375) (2.34,-0.641) (2.36,-0.826) (2.38,-0.702) (2.40,-0.573) (2.42,-0.690) (2.44,-0.552) (2.46,-0.426) (2.48,-0.774) (2.50,-0.821) (2.52,-0.722) (2.54,-0.822) (2.56,-1.233) (2.58,-1.265) (2.60,-1.090) (2.62,-0.752) (2.64,-1.175) (2.66,-0.629) (2.68,-0.842) (2.70,-0.633) (2.72,-0.374) (2.74,-0.170) (2.76,-0.136) (2.78,-0.366) (2.80,-0.233) (2.82,-0.372) (2.84,-0.869) (2.86,-0.286) (2.88,-0.138) (2.90,0.228) (2.92,0.041) (2.94,0.334) (2.96,0.645) (2.98,0.943) (3.00,0.920) (3.02,0.661) (3.04,0.612) (3.06,0.160) (3.08,-0.184) (3.10,-0.162) (3.12,-0.357) (3.14,-0.383) (3.16,-0.397) (3.18,-0.003) (3.20,0.257) (3.22,0.243) (3.24,0.480) (3.26,0.314) (3.28,0.240) (3.30,0.418) (3.32,0.161) (3.34,0.102) (3.36,-0.164) (3.38,-0.137) (3.40,0.192) (3.42,0.038) (3.44,0.082) (3.46,-0.099) (3.48,-0.327) (3.50,-0.568) (3.52,-0.265) (3.54,0.212) (3.56,0.305) (3.58,0.155) (3.60,0.293) (3.62,0.221) (3.64,0.376) (3.66,0.427) (3.68,0.475) (3.70,0.586) (3.72,0.459) (3.74,0.369) (3.76,0.017) (3.78,-0.072) (3.80,-0.355) (3.82,-0.376) (3.84,-0.560) (3.86,-0.529) (3.88,-0.428) (3.90,-0.703) (3.92,-0.852) (3.94,-1.023) (3.96,-0.854) (3.98,-0.495) (4.00,-0.315) (4.02,-0.379) (4.04,-0.078) (4.06,-0.379) (4.08,-0.072) (4.10,-0.202) (4.12,-0.754) (4.14,-0.208) (4.16,0.111) (4.18,-0.097) (4.20,-0.062) (4.22,-0.103) (4.24,-0.081) (4.26,-0.376) (4.28,-0.505) (4.30,-0.401) (4.32,-0.348) (4.34,-0.102) (4.36,-0.084) (4.38,0.388) (4.40,0.237) (4.42,0.284) (4.44,-0.105) (4.46,-0.353) (4.48,-0.081) (4.50,-0.267) (4.52,-0.156) (4.54,-0.165) (4.56,-0.371) (4.58,-0.433) (4.60,-0.527) (4.62,-0.614) (4.64,-0.451) (4.66,-0.312) (4.68,-0.421) (4.70,-0.594) (4.72,-0.526) (4.74,-0.535) (4.76,-0.319) (4.78,0.219) (4.80,0.383) (4.82,0.368) (4.84,0.327) (4.86,0.149) (4.88,-0.036) (4.90,-0.236) (4.92,-0.310) (4.94,-0.389) (4.96,-0.233) (4.98,-0.308) (5.00,-0.341)};
    \end{tikzpicture}
    \caption{Left: three paths of Brownian motion, simulated with $h = 0.02$. Right: three paths of the Ornstein--Uhlenbeck process $dX = -X\,dt + \sqrt{2}\,dB$, started at $2.5$, $-2$ and $0$, together with the solutions $X(0)e^{-t}$ of the ODE without noise (dashed). The process forgets its starting point and fluctuates around $0$ with the stationary distribution $\mathcal{N}(0, 1)$; about $95\%$ of the time, it stays between $-2$ and $2$.}
    \label{L5-fig:sde}
\end{figure}

\section{Langevin Dynamics and Diffusion Models}\label{L5-sec:langevin}


The Ornstein--Uhlenbeck process is an example of a general principle: a gradient flow with the right amount of noise samples from a probability distribution.

\begin{theorem}[Langevin dynamics]\label{L5-thm:langevin}
    Let $U: \mathbb{R}^d \to \mathbb{R}$ be smooth with $\int e^{-U(\mathbf{x})}\, d\mathbf{x} < \infty$. Then the SDE $d\mathbf{X} = -\nabla U(\mathbf{X})\,dt + \sqrt{2}\,d\mathbf{B}$, where $\mathbf{B}$ is a vector of $d$ independent Brownian motions, has the stationary density $p(\mathbf{x}) \propto e^{-U(\mathbf{x})}$, and under mild conditions, the distribution of $\mathbf{X}(t)$ converges to it.
\end{theorem}

We do not prove the convergence, but we can see why $e^{-U}$ is stationary. Without noise, the density $p_t$ of $\mathbf{X}(t)$ under the flow of the vector field $-\nabla U$ satisfies the continuity equation $\frac{\partial p_t}{\partial t} = -\Div(-p_t\nabla U)$ (Proposition~\ref{L5-prop:continuity}). The noise adds a diffusion term, and one can show that $p_t$ satisfies the Fokker--Planck equation
\[\frac{\partial p_t}{\partial t} = \Div(p_t\nabla U) + \Delta p_t = \Div(p_t\nabla U + \nabla p_t),\]
with the Laplacian $\Delta p_t = \Div\nabla p_t$ (Lecture~3). For $p \propto e^{-U}$, we have $\nabla p = -p\nabla U$, so the right side vanishes, and $p$ does not change in time. For $U = 0$, the Fokker--Planck equation is the heat equation $\frac{\partial p_t}{\partial t} = \Delta p_t$, which describes how the density of $\sqrt{2}\,\mathbf{B}(t)$ spreads out. For $U(x) = \frac{x^2}{2s^2}$, the SDE is the Ornstein--Uhlenbeck process with $\theta = \frac{1}{s^2}$ and $\sigma = \sqrt{2}$, whose stationary variance $\frac{\sigma^2}{2\theta} = s^2$ agrees with the density $e^{-U} \propto e^{-x^2/(2s^2)}$ of $\mathcal{N}(0, s^2)$. The Euler--Maruyama scheme gives the unadjusted Langevin algorithm
\[\mathbf{x}_{k+1} = \mathbf{x}_k - h\nabla U(\mathbf{x}_k) + \sqrt{2h}\,\mathbf{z}_k, \qquad \mathbf{z}_k \sim \mathcal{N}(\mathbf{0}, I),\]
which is gradient descent on $U$ plus Gaussian noise. To sample from a posterior $p(\boldsymbol{\theta} \mid \mathcal{D}) \propto e^{-U(\boldsymbol{\theta})}$, we take $U(\boldsymbol{\theta}) = -\log p(\mathcal{D} \mid \boldsymbol{\theta}) - \log p(\boldsymbol{\theta})$, the training loss plus the regularizer. Without the noise, the iteration converges to the MAP estimate; with the noise, it explores the posterior. The discretization introduces a small bias: for the target $\mathcal{N}(2, 0.25)$ and $h = 0.01$, the stationary variance of the algorithm is $0.2551$ instead of $0.25$, by Example~\ref{L5-ex:ou}. The Metropolis-adjusted Langevin algorithm removes this bias by accepting or rejecting each step as in Theorem~\ref{L4-thm:metropolis}, and stochastic gradient Langevin dynamics uses mini-batch gradients to make the method feasible for large models.

\begin{remark}[SDEs in machine learning]
    \begin{enumerate}
        \item Diffusion models for image generation use an SDE of Ornstein--Uhlenbeck type that gradually turns images into Gaussian noise. The time-reversed process is again an SDE, whose drift involves the score $\nabla\log p_t(\mathbf{x})$ of the noised data distribution at time $t$. A neural network learns this score, with a variant of the score matching loss of Lecture~3, and new images are generated by simulating the reverse SDE, starting from pure noise.
        \item The noise of mini-batch gradients makes SGD with a small step size behave approximately like the Euler--Maruyama scheme for an SDE $d\boldsymbol{\theta} = -\nabla L(\boldsymbol{\theta})\,dt + \sqrt{\eta\,\Sigma(\boldsymbol{\theta})}\,d\mathbf{B}$, where $\Sigma$ is the covariance matrix of the mini-batch gradients. This heuristic is used to study which minima SGD prefers.
    \end{enumerate}
\end{remark}

\section{Key Takeaways}
\begin{itemize}[leftmargin=*]
    \item An ODE $\mathbf{x}' = F(\mathbf{x})$ with a Lipschitz continuous $F$ has a unique solution for each initial value. Linear systems $\mathbf{x}' = A\mathbf{x}$ are solved with eigenvectors, $\mathbf{x}(t) = \sum_i c_ie^{\lambda_i t}\mathbf{v}_i = e^{At}\mathbf{x}_0$, and the signs of the real parts of the eigenvalues decide stability. Along the gradient flow $\mathbf{x}' = -\nabla f(\mathbf{x})$, $f$ decreases; gradient descent is Euler's method for it, and its step size limit $\frac{2}{\lambda_{\max}}$ is the stability limit of Euler's method for stiff problems.
    \item The flow $\varphi_t$ of a vector field $F$ expands volumes at the rate $\Div F$; for linear fields, $\det e^{At} = e^{t\tr A}$. A probability density moved by the flow satisfies the continuity equation $\frac{\partial p_t}{\partial t} = -\Div(p_tF)$, and along each solution, $\frac{d}{dt}\log p_t(\mathbf{x}(t)) = -\Div F(\mathbf{x}(t))$, which continuous normalizing flows use instead of a Jacobian determinant.
    \item Brownian motion has independent Gaussian increments $B(t) - B(s) \sim \mathcal{N}(0, t - s)$ and continuous, nowhere differentiable paths. An SDE $dX = \mu(X)\,dt + \sigma(X)\,dB$ is simulated by Euler--Maruyama, $X_{k+1} = X_k + \mu(X_k)h + \sigma(X_k)\sqrt{h}\,Z_k$: noise of size $\sqrt{h}$, not $h$. Langevin dynamics $d\mathbf{X} = -\nabla U\,dt + \sqrt{2}\,d\mathbf{B}$ samples from $p \propto e^{-U}$, the stationary solution of the Fokker--Planck equation; its discretization is gradient descent plus noise, and diffusion models generate data by reversing a noising SDE.
\end{itemize}

\section{Exercises}\label{L5-sec:exercises}
\begin{enumerate}[leftmargin=*]
    \item Solve the ODE $\mathbf{x}' = A\mathbf{x}$ with $A = \begin{pmatrix} 1 & 2 \\ 2 & 1 \end{pmatrix}$ and $\mathbf{x}(0) = (1, 0)^\top$, using the eigenvectors of $A$. Which initial values give solutions that stay bounded as $t \to \infty$? Compute $e^{At}$.
    \item\label{L5-ex5:euler-skew} \begin{enumerate}[label=(\alph*)]
        \item Show that Euler's method for $x' = -\lambda x$ with $\lambda > 0$ gives $x_k = (1 - h\lambda)^k x_0$. For which step sizes $h$ does the numerical solution tend to $0$, like the exact solution? Relate the answer to Theorem~\ref{L3-thm:gd-quadratic}.
        \item Let $A$ be skew-symmetric. Show that $Q(t) = e^{At}$ is orthogonal for all $t$, by differentiating $Q(t)^\top Q(t)$. Show that Euler's method for $\mathbf{x}' = A\mathbf{x}$ increases $\norm{\mathbf{x}_k}$ in every step (unless $A\mathbf{x}_k = \mathbf{0}$).
    \end{enumerate}
    \item\label{L5-ex5:nilpotent} Let $A = \begin{pmatrix} -1 & 2 \\ 0 & -1 \end{pmatrix} = -I + N$ with $N = \begin{pmatrix} 0 & 2 \\ 0 & 0 \end{pmatrix}$.
    \begin{enumerate}[label=(\alph*)]
        \item Show that $A$ is not diagonalizable, and that $N^2 = 0$.
        \item Show that $e^{At} = e^{-t}(I + tN)$, by checking that it equals $I$ at $t = 0$ and that its derivative is $A$ times itself. Solve $\mathbf{x}' = A\mathbf{x}$ with $\mathbf{x}(0) = (0, 1)^\top$.
        \item Check Liouville's formula $\det e^{At} = e^{t\tr A}$ for this matrix. By which factor does the flow change areas in the time $t$?
    \end{enumerate}
    \item Let $\mathbf{x}' = A\mathbf{x}$, and let $\mathbf{x}(0)$ be random with the density $p_0$. Use the change of variables formula (Lecture~3, Remark~\ref{L3-rem:change-of-variables}) and Proposition~\ref{L5-prop:liouville} to show that $\log p_t(\mathbf{x}(t)) = \log p_0(\mathbf{x}(0)) - t\tr A$, and check that this agrees with Proposition~\ref{L5-prop:continuity}. What is the distribution of $\mathbf{x}(t)$ if $\mathbf{x}(0) \sim \mathcal{N}(\mathbf{0}, I)$?
    \item (Brownian motion and the Ornstein--Uhlenbeck process.)
    \begin{enumerate}[label=(\alph*)]
        \item Compute $\Var(B(2) - B(0.5))$, $\Cov(B(1), B(3))$ and $\Corr(B(1), B(3))$.
        \item Let $c > 0$. Show that $W(t) = \frac{1}{\sqrt{c}}B(ct)$ has the same mean and covariance functions as $B(t)$. (In fact, $W$ is again a Brownian motion: zooming into a Brownian path by the factor $c$ in time and $\sqrt{c}$ in space gives a Brownian path.)
        \item Let $X_{k+1} = (1 - \theta h)X_k + \sigma\sqrt{h}\,Z_k$ with $X_0 = x_0$ and i.i.d.\ $Z_k \sim \mathcal{N}(0, 1)$. Show that $\E[X_k] = (1 - \theta h)^kx_0$ and $\Var(X_k) = v_h\big(1 - (1 - \theta h)^{2k}\big)$, with $v_h$ from Example~\ref{L5-ex:ou}. What happens for $\theta h > 2$?
    \end{enumerate}
\end{enumerate}

\section{Programming Exercises}\label{L5-sec:programming}

% Inline code cannot be hyphenated, so allow a little extra stretch in this section.
\setlength{\emergencystretch}{3em}

As in the previous lectures, implement the methods yourself with NumPy, and use library routines (such as \py{scipy.integrate.solve_ivp}) only to check your results. Start each file with
\begin{lstlisting}
import numpy as np
import matplotlib.pyplot as plt
rng = np.random.default_rng(0)
\end{lstlisting}

\begin{enumerate}[label=\textbf{P\arabic*.}, leftmargin=*]
    \item \textbf{ODEs and gradient flow.}
    \begin{itemize}
        \item Implement Euler's method for $\mathbf{x}' = F(\mathbf{x})$. For $A = \begin{pmatrix} -2 & 1 \\ 1 & -2 \end{pmatrix}$ and $\mathbf{x}_0 = (2, 0)^\top$, compare the numerical solution at $t = 1$ with the exact solution of Example~\ref{L5-ex:linear-ode}, for $h = 0.1, 0.05, 0.025, \ldots$. Plot the error against $h$ on a log-log scale; what is the slope?
        \item Implement the fourth-order Runge--Kutta method, $\mathbf{x}_{k+1} = \mathbf{x}_k + \frac{h}{6}(\mathbf{k}_1 + 2\mathbf{k}_2 + 2\mathbf{k}_3 + \mathbf{k}_4)$ with $\mathbf{k}_1 = F(\mathbf{x}_k)$, $\mathbf{k}_2 = F(\mathbf{x}_k + \frac{h}{2}\mathbf{k}_1)$, $\mathbf{k}_3 = F(\mathbf{x}_k + \frac{h}{2}\mathbf{k}_2)$, $\mathbf{k}_4 = F(\mathbf{x}_k + h\mathbf{k}_3)$, and repeat the error plot.
        \item Plot the vector field of Figure~\ref{L5-fig:ode} with \py{plt.quiver}, together with numerical solutions. Do the same for the gradient flow of the Rosenbrock function, and compare the solution of the flow (computed with a small step) with the gradient descent iterates of Programming Exercise P8 of Lecture~3 for $\eta = 10^{-3}$ and $\eta = 3 \cdot 10^{-3}$.
    \end{itemize}

    \item \textbf{Flows and volumes.}
    \begin{itemize}
        \item Draw $10\,000$ samples from $\mathcal{N}(\mathbf{0}, I)$ in $\mathbb{R}^2$, and move them with the flow of $\mathbf{x}' = A\mathbf{x}$ for the matrix $A$ of Figure~\ref{L5-fig:ode} up to $t = 1$, using your Runge--Kutta method from P1. Compare the sample covariance matrix with $e^{At}e^{A^\top t}$, and its determinant with $e^{2t\tr A}$.
        \item The pendulum $x_1' = x_2$, $x_2' = -\sin x_1$ has $\Div F = 0$, and with friction, $x_2' = -\sin x_1 - \gamma x_2$, it has $\Div F = -\gamma$. Represent a small disk around $(1, 0)$ by $200$ points on its boundary circle, move them with the flow, and compute the enclosed area over time with the shoelace formula $\frac{1}{2}\sum_i (x_iy_{i+1} - x_{i+1}y_i)$. Compare with the prediction of the divergence for $\gamma = 0$ and $\gamma = 0.5$, and plot the deformed disks.
    \end{itemize}

    \item \textbf{Stochastic differential equations.}
    \begin{itemize}
        \item Simulate $10\,000$ Brownian paths on $[0, 1]$ with the step $h = 0.001$. Check $\Var(B(t)) = t$ and $\Cov(B(0.3), B(0.8)) = 0.3$ empirically, and plot a few paths. Zoom into one path: plot it on $[0.5, 0.5 + \delta]$ for $\delta = 0.1, 0.01, 0.001$ (simulating the path with a step much smaller than $\delta$). Why does it never look smooth?
        \item Simulate the Ornstein--Uhlenbeck process $dX = -\theta X\,dt + \sigma\,dB$ with Euler--Maruyama for $\theta = 1$, $\sigma = \sqrt{2}$, starting $10\,000$ paths at $x_0 = 3$. Plot the mean and the variance of $X(t)$ against $t$, together with $x_0e^{-t}$ and $1 - e^{-2t}$. Compare the histogram at $t = 5$ with $\mathcal{N}(0, 1)$, and check $v_h$ from Example~\ref{L5-ex:ou} for $h = 0.5$.
        \item (Langevin sampling.) Use the unadjusted Langevin algorithm to sample from the mixture $p = \frac{1}{2}\mathcal{N}(-2, 0.5^2) + \frac{1}{2}\mathcal{N}(2, 0.5^2)$, i.e., $U = -\log p$ (compute $\nabla U$ by hand). Compare a histogram of one long chain with $p$, and count how often the chain moves between the two modes. Then sample the posterior of the logistic regression model of Lecture~3 on a small two-dimensional data set, with the prior $\mathcal{N}(\mathbf{0}, I)$, and plot the samples of $\mathbf{w}$ together with the MAP estimate.
    \end{itemize}

\end{enumerate}
\endgroup
